\chapter{Decay, Confinement, and Symmetry Breaking}
\label{chap:decay-confinement-symmetry}

This chapter is the particle-physics frontier, and it closes Part IV. Here, more than anywhere in the book, the
physics outruns the finite model: most of these eight effects rest on the name-only floor, their physics
written out in full and the model claiming only the name. But the frontier keeps its walls and its counts as
well --- a decay that removes a branch, a conservation that balances a split, a winding that must be an integer,
and a vacuum that pulls by a shadow. And it is, for all the modesty of its floors, among the math-richest
chapters in the book: confinement's growing potential, the sombrero that breaks a symmetry, the non-abelian
field strength, the Casimir mode-count --- each is written out, every object given its picture, even where the
model claims only the label beneath it. All four honest floors appear here, in this one chapter.

\section{The Alpha Decay Effect: a branch removed}
\label{se:alpha-decay}

Alpha decay is usually told as tunneling --- a particle leaking through a barrier it has not the energy to
cross. Read off finite records, no tunneling is assumed. An unstable nucleus carries an irreconcilable overlap,
a branch that cannot be made consistent --- here, an alpha particle, a tightly bound clump of two protons and
two neutrons that the nucleus can no longer hold in one consistent record --- and decay is the unique minimal
boundary update that removes it: the inconsistent branch is eliminated and the record evolves on the one that
remains. The removal is not timed but only probable per unit time, each nucleus carrying the same fixed chance
of the repair in any interval. Summed over many nuclei, that constant per-capita rate gives the exponential
law,
\begin{equation}
  N(t) = N_0\, e^{-\lambda t},
\end{equation}
the population falling by a fixed fraction in each span, with a half-life $t_{1/2} = \ln 2 / \lambda$ at which
half remain. The honest floor is a no-go --- the inconsistent branch cannot be kept, and its removal is forced.
The reading is that alpha decay is the discarding of a branch, not the leaking of a particle. No hidden force
and no barrier-crossing is needed: the nucleus carries a branch it cannot reconcile, the minimal repair removes
it, and the exponential law is the statistics of that forced removal --- a wall the record meets, not a barrier
it sneaks through.

\section{The Gamma Decay Effect: a conservative split}
\label{se:gamma-decay}

Gamma decay is a reversible, conservative split. A nucleus left in an excited state --- often by an earlier
alpha or beta decay --- carries unresolved internal curvature and relaxes to its ground state by emitting a
photon that carries off exactly the difference,
\begin{equation}
  E_{\text{excited}} = E_{\text{ground}} + E_\gamma .
\end{equation}
The nucleus, like the atom, has discrete energy levels, and the gamma photon's energy is the gap between two of
them --- which is why gamma spectra are sharp lines, each line a particular transition, the nuclear fingerprint
by which a decaying isotope is known. Nothing is lost: the excited count equals the ground count plus the
photon, and the split balances exactly. The honest floor is a finite count obligation --- the conservation is a
count, the excited total equal to the ground total plus the emitted photon. The reading is that gamma decay is
conservation made into a transition. The excited nucleus does not lose its energy; it partitions it, handing
the photon exactly the excess so that the books balance across the split, the sharp line the signature of an
exact count.

\section{The Neutrino Effect: the messenger that arrives first}
\label{se:neutrino}

The neutrino barely interacts, and that is the whole of the effect. It carries no charge and feels neither
electromagnetism nor the strong force; only the weak interaction touches it, and the weak interaction is
feeble, so a neutrino passes through matter almost as though it were not there --- light-years of lead would
scarcely stop it. Because almost nothing slows it, a neutrino travels nearly unimpeded and saturates the speed
limit $c$. The proof arrived from the sky in 1987: when a star exploded as supernova SN1987A, its neutrinos
reached Earth about three hours before its light, because the neutrinos streamed straight out of the collapsing
core while the photons scattered their slow way out through the star's shattering envelope. The messenger that
interacts least arrives first.

The honest floor is name-only --- the physics of the weakly-interacting messenger is written as physics, and
the finite model claims only the name. The reading is that the neutrino is named, its early arrival is the
physics of a particle that almost nothing slows, and the finite model is the label beneath it: the messenger
that arrives first because the medium has almost no hold on it.

\section{The Strong Interaction Effect: strain self-confines}
\label{se:strong-interaction}

The strong interaction confines, and it does so by a force unlike any other: one that does not weaken with
distance but strengthens. Where a weak disturbance disperses through the record as a linear smoothing, the
strong regime is nonlinear --- the strain does not spread but self-confines, and the combined strain of two
charges is not the sum of their separate strains,
\begin{equation}
  \Sigma(q_i, q_j) \neq \Sigma(q_i) + \Sigma(q_j) .
\end{equation}
The potential between two quarks grows with their separation rather than falling off,
\begin{equation}
  V(r) \sim \sigma r,
\end{equation}
a straight line of slope $\sigma$, the string tension: the colour field between the quarks collapses into a
narrow flux tube of fixed tension, and pulling the quarks apart lengthens the tube at a constant energy cost
per unit length. Pull hard enough and it is cheaper for the tube to snap --- the stored energy making a fresh
quark--antiquark pair at the break --- than to stretch further, so a quark is never isolated; it comes away
clothed in a new partner, and no free quark is ever seen. This is visible in every high-energy collision: a
quark knocked hard out of a proton does not fly off alone but dresses itself, on the way out, into a
collimated spray of ordinary particles --- a jet --- the snapping flux tube made into a shower; detectors see
jets, never quarks. The same force has an opposite face at short range:
asymptotic freedom. The coupling runs weak at small separation, so quarks squeezed close together inside a
hadron behave almost as free particles, while the coupling grows strong at large separation and confines them.
Weak when near, unbreakable when far --- the reverse of every other force.

The honest floor is name-only --- the confinement physics is written in full, and the finite model claims only
the name. The reading is that confinement is named, the linear potential $V \sim \sigma r$, the asymptotic
freedom, and the super-additive strain are the physics, and the finite model is the label: a force that grows
with distance, so its charges are never isolated.

\section{The Sombrero Potential: symmetry breaks, mass appears}
\label{se:sombrero}

Spontaneous symmetry breaking is the sombrero potential. A complex scalar field $\phi$ sits in a potential
shaped like a Mexican hat,
\begin{equation}
  V(\phi) = \lambda\bigl(|\phi|^2 - v^2\bigr)^2,
\end{equation}
symmetric about the center but with its minimum not at the center --- where the potential has a bump --- but
around a whole circle of radius $|\phi| = v$, the brim of the hat. The field must come to rest somewhere on
that circle, and in choosing a point it breaks the symmetry: the law is symmetric under rotations of $\phi$,
but the chosen ground state is not, singling out one direction around the brim. Write the field as a
displacement from the chosen point,
\begin{equation}
  \phi = (v + h)\,e^{i\theta},
\end{equation}
and the two excitations split in character. The radial one, $h$, climbs the wall of the hat and costs energy
--- it is massive, the Higgs mode. The angular one, $\theta$, runs around the flat brim at no cost --- it is
massless, a Goldstone mode. And here is the mechanism: when the broken symmetry is a gauge symmetry, the
massless Goldstone $\theta$ is not a physical particle at all but a phase the gauge field $A_\mu$ can absorb.
The gauge boson swallows the Goldstone, gaining a longitudinal mode and, with it, a mass --- a massless
force-carrier eats the massless angular excitation and becomes a massive one, which is how the weak bosons get
their mass and the force its short range. The radial mode is not a fiction either: the Higgs boson, the
quantum of the field's climb up the wall of the hat, was found at the Large Hadron Collider in 2012 at a mass
near 125~GeV --- the last piece of the Standard Model seen at last.

The honest floor is name-only --- the symmetry-breaking and mass-generation physics is written in full, and the
finite model claims only the name. The reading is that the Higgs mechanism is named, the sombrero $V(\phi)$, the
broken circle, and the absorbed Goldstone are the physics, and the finite model is the label: a symmetry the
field breaks by choosing, and a mass the gauge boson gains by swallowing what is left.

\section{The Yang--Mills Effect: the gauge group}
\label{se:yang-mills}

The Standard Model's forces are a gauge symmetry. Each class of label defines a distinct mode of refinement,
and the requirement that local corrections compose consistently forces each into a compact local symmetry;
together they are the Standard-Model group,
\begin{equation}
  U(1) \times SU(2) \times SU(3),
\end{equation}
electromagnetism, the weak force, and the strong force as three local symmetries. What appears in the smooth
shadow as a gauge field $A_\mu$ is, beneath it, the connection read off how local frames must be adjusted to
compose consistently --- the same connection whose curvature the Dirac chapter wrote as $[D_\mu, D_\nu]$. But
here the connection has a feature electromagnetism lacks. For the abelian $U(1)$ the field strength is the
plain curl,
\begin{equation}
  F_{\mu\nu} = \partial_\mu A_\nu - \partial_\nu A_\mu,
\end{equation}
and the photon carries no charge of its own. For the non-abelian $SU(2)$ and $SU(3)$ the gauge fields do not
commute, and the field strength gains a term built from their commutator,
\begin{equation}
  F_{\mu\nu} = \partial_\mu A_\nu - \partial_\nu A_\mu + ig\,[A_\mu, A_\nu],
\end{equation}
the extra piece a self-interaction: the gauge fields carry the very charge they mediate, so gluons pull on
gluons and the weak bosons on one another, where photons pass through each other untouched. That self-coupling
is the root of confinement and of the whole nonlinearity of the strong and weak forces.

The honest floor is name-only --- the gauge structure is written as physics, and the finite model claims only
the name. The reading is that the Standard-Model gauge group is named, the product $U(1) \times SU(2) \times
SU(3)$, its connection $A_\mu$, and the non-abelian field strength are the physics, and the finite model is the
label: three forces as three compact symmetries, the gauge field their smooth shadow, self-interacting where it
does not commute.

\section{The Casimir Effect: the vacuum pulls by a shadow}
\label{se:casimir}

Two uncharged plates set close together attract, and the pull comes from the vacuum. The vacuum is not empty:
every mode of the electromagnetic field carries a zero-point energy $\tfrac12\hbar\omega$ --- the same restless
half-quantum the oscillator could not give up in the uncertainty chapter --- and the vacuum's energy is the sum
over all modes,
\begin{equation}
  \sum \tfrac12\hbar\omega,
\end{equation}
a sum that is enormous, and the same outside the plates as in, until the plates restrict it. Between two
conducting plates a distance $d$ apart, only modes whose wavelengths fit --- a whole number of half-waves
across the gap --- are allowed, and the long-wavelength modes that do not fit are excluded. So the interior
carries fewer modes than the exterior, and the zero-point sum inside falls below the sum outside. The imbalance
is a pressure: the fuller vacuum outside presses the plates together harder than the thinner vacuum inside
presses them apart. Subtract the two infinite sums, with care, and the finite difference is an attractive force
per unit area,
\begin{equation}
  \frac{F}{A} = -\frac{\pi^2 \hbar c}{240\, d^4},
\end{equation}
growing sharply, as the inverse fourth power, as the plates near --- a force since measured, and matched.

The honest floor is a smooth-shadow analogy --- the discrete interior mode count is the model, the continuum
Casimir pressure its smooth shadow, and the bridge from the missing modes to the $1/d^4$ law is an analogy,
named as one. The reading is that the Casimir force is a counted absence. The plates attract because there are
fewer modes within than without, and the continuum law $F/A \propto 1/d^4$ is the smooth shadow of that finite
deficit --- the vacuum pulling by what it is not allowed to carry, the oscillator's $\tfrac12\hbar\omega$ now
summed over every mode and felt as a force.

\section{The Topological Integer Count: quantization by winding}
\label{se:topological}

Some quantities are quantized not by a spectrum but by a topology --- not because a count fills a ladder, but
because a configuration wraps. A field whose phase $\theta$ turns as one goes around a loop carries a winding
number, the total turning divided by a full turn,
\begin{equation}
  w = \frac{1}{2\pi}\oint d\theta \;\in\; \mathbb{Z},
\end{equation}
and it is necessarily an integer: the phase must return to itself after the circuit, so the total turning is a
whole number of full cycles, never a fraction. This is the Chern number, and it is what the quantum Hall effect
measures --- the Hall conductance locked to integer multiples of a fundamental unit, flat plateaus that do not
drift with the sample's details because an integer cannot drift, the TKNN integer counting how the state winds
over the space of momenta. No fractional thread is admissible: the record either contains a whole thread or it
does not, and a half-turn that fails to close is not a configuration at all.

The honest floor is a finite count obligation --- the winding is an integer count, and fractional values are
inadmissible. The reading is that topological quantization is counting by winding. The integer is not the rung
of a ladder but the number of times a configuration wraps, and because a thread either closes or does not, the
count is exact and whole --- quantization with no continuum in sight, the conductance pinned to an integer by
the impossibility of a fractional turn.

\bigskip

With the frontier, Part IV closes, and with it the longest movement of the book. All four honest floors met in
this one chapter: a no-go in the discarded branch of alpha decay, a finite count in the balanced split of gamma
decay and the integer winding, a smooth shadow in the Casimir deficit, and a name-only label beneath the
neutrino, the strong force, the Higgs, and the gauge group --- the instrument's full vocabulary of honesty,
gathered at the frontier where the physics runs furthest ahead. The quantum began as a doorway --- the count
carried as an amplitude, diffusion run in imaginary time --- and through the part the amplitude became a
holonomy that a loop charges to $\pm 1$, a particle that a detector proves by coincidence, and, at the
frontier, a name the physics still outruns. The count became an amplitude, the amplitude a loop charge, the
loop charge a particle proved at a detector; and where the model could prove no more, it named honestly and
wrote the physics in full. Part V turns next to relativity and gravitation, where the count acquires a geometry
and the record's own measure bends.
